Sentiment scores for text are usually reported as one bounded number per document, such as a class probability or a normalised sum of word values. Such numbers cannot be recombined: the score of a chapter does not follow from the scores of its sentences, so long texts are cut into chunks and averaged, and the result depends on where the cuts fall. Here a measure of sentiment is introduced, the felicific signature, whose value for a span of text is computed exactly from the values of its parts when the text is divided between sentences. Each word carries positive charge (hedons) and negative charge (dolors) taken from human word ratings, and a span is summarised by 19 numbers: counts, position-weighted sums of charge, and the charge under negation and other operators. One pass over a book therefore yields the signature of every sentence, every chapter and the whole. Among position weights that admit such a merge, polynomials are the only continuous functions of relative position. On five corpora with continuous human ratings the calibrated score correlates with the ratings at $r=\meanSgw$, comparable to VADER (\meanVader), the strongest of the three other lexicon methods tested, and below two off-the-shelf transformer classifiers (\meanDistil{} and \meanRoberta), at \speedTimes{} times their speed on one processor core. On whole film reviews a simpler lexicon method agrees better with the critics' ratings, and wherever agreement above the sentence could be checked, averaged sentence scores from another method tracked the human ratings better than the merged signature. In \vaderBooks{} novels a conventional chunked score moves from \vaderOne{} to \vaderChapter{} as the chunks grow, whereas merged signatures do not change. The hypothesis that sentiment is front-loaded is not supported. Words carrying no sentiment dilute a rating: the planned test supports this in \winsNoDil{} of six corpora, and in exploratory analyses of human phrase ratings the effect appears from either side and a part's weight grows with its share of the words. Compared with student essays on the same assignments, essays written by four language models close on a more positive note, as predicted, and, in an exploratory comparison, place less of their charge after intensifiers. The signature gives the same result however a text is divided between sentences, at a small fraction of a transformer's cost; on single texts, transformer classifiers remain more accurate.
